\chapter{Arbitrage d'indice et arbitrage de panier}\label{ch:hf:index-and-basket-arbitrage}

Le contrat à terme sur indice monte, mais cinq cents actions n'ont pas encore bougé. Acheter les cinq cents prend plus de temps que l'écart ne dure, alors l'arbitragiste
achète les trente qui comptent et accepte le risque sur les autres. Dans l'indice simulé de cinquante noms de ce chapitre, un arbitragiste qui envoie une jambe
toutes les vingt microsecondes gagne le plus par transaction avec un panier de cinq noms, et le meilleur ratio d'edge sur risque avec quinze~; celui qui envoie une jambe
toutes les deux microsecondes gagne le plus avec dix, et son meilleur ratio d'edge sur risque correspond à l'ensemble des cinquante.

\section{La base futures--comptant à haute fréquence}

L'arbitrage d'indice (livre~9, chapitre~26) maintient le contrat à terme à son \omterm{def:hf:fair-value:fair}{prix juste} (livre~1, chapitre~21), $F=S\,e^{(r-q)T}$, dans une bande fixée par les
coûts de négociation du contrat à terme et du panier. À l'horizon de quelques jours, c'est un trade de financement~; à l'horizon de quelques millisecondes, c'est un trade lead--lag
(chapitre~8)~: le contrat à terme, un instrument doté d'un seul carnet d'ordres, réagit aux nouvelles qui touchent tout le marché avant que les centaines de carnets d'ordres
des composantes de l'indice ne le fassent, et pendant un instant la base sort de sa bande parce que le côté comptant est en retard.

La largeur de la bande dépend de la liquidité. Roll, Schwartz et Subrahmanyam ont étudié la base du contrat à terme sur l'indice NYSE Composite par rapport à la
liquidité agrégée du NYSE sur quelque 3\,000 jours de séance~: la liquidité et la valeur absolue de la base se prévoient mutuellement, avec une causalité de Granger
bilatérale à court horizon, et les chocs sur la base prédisent une liquidité ultérieure. L'arbitrage maintient les prix ensemble lorsque le marché est assez liquide pour le permettre.

\section{Paniers partiels et risque de réplication}

\begin{definition}[Panier partiel]\label{def:hf:index-and-basket-arbitrage:partial}
Un \emph{panier partiel}\index{panier partiel} est un sous-ensemble des composantes d'un indice, pondérées de manière à suivre l'indice aussi fidèlement que possible,
négocié à la place du panier complet afin de gagner du temps et de réduire les coûts, au prix d'une erreur de réplication.
\end{definition}

Avec une matrice de covariance $\Sigma$ des rendements des composantes et des pondérations de l'indice $w$, la variance de réplication d'un panier $h$ vaut
$(h-w)^\top\Sigma(h-w)$. Pour un ensemble donné $S$ de noms, les poids minimisants résolvent $\Sigma_{SS}h_S=(\Sigma w)_S$~; l'ensemble est choisi de façon gloutonne, un nom
à la fois, chacun étant celui qui réduit le plus la variance de réplication
(\cref{lst:hf:index-and-basket-arbitrage:partial}). Dans un modèle à un facteur, les premiers noms retenus sont les plus gros et les plus représentatifs~; le bêta du panier
se rapproche de un à mesure que l'on ajoute des noms. Dans l'indice de cinquante noms du chapitre, les dix plus grandes capitalisations pèsent la moitié de son poids~; un panier de cinq noms détient
73~\% du notionnel avec un bêta de 0.83, et un panier de quinze noms 90~\% avec un bêta de 0.94.

\section{Risque de jambes et risque d'exécution}

\begin{definition}[Risque d'exécution par jambes]\label{def:hf:index-and-basket-arbitrage:legging}
Le \emph{risque d'exécution par jambes}\index{risque d'exécution par jambes} (legging risk) est le risque que les prix des jambes d'une transaction à plusieurs instruments évoluent entre l'exécution de la
première jambe et celle de la dernière, de sorte que la transaction s'achève à des prix moins favorables que ceux qui la justifiaient, ou ne s'achève pas.
\end{definition}

L'edge de l'arbitragiste sur une jambe est le mouvement qu'il capte en achetant avant que les teneurs de marché de cette action ne mettent à jour leurs cotations. Les jambes partent
l'une après l'autre par la passerelle d'ordres de la société~: avec $k$ jambes envoyées toutes les $\delta$ microsecondes, la dernière arrive $k\delta$ après la première. Si la cotation d'une action
se met à jour après un délai exponentiel de moyenne $\tau$, la jambe $j$ arrive à temps avec la probabilité $e^{-j\delta/\tau}$. Davantage de noms réduisent
l'erreur de réplication et pèsent davantage dans la transaction, mais les jambes tardives arrivent après que leurs actions ont bougé, et chaque nom ajouté est plus petit et a un
spread plus large. La taille optimale du panier est celle qui équilibre ces trois effets.

\begin{omfigure}
\begin{tikzpicture}
\begin{axis}[width=11cm,height=5.4cm,xlabel={\small noms dans le panier},ylabel={\small points de base par transaction},xmode=log,log basis x=10,
  xtick={1,2,5,10,20,50},xticklabels={1,2,5,10,20,50},xmin=0.9,xmax=55,ymin=0,ymax=1.8,tick label style={font=\footnotesize},grid=major,
  grid style={gray!25},table/col sep=comma,legend style={font=\scriptsize,at={(0.5,-0.3)},anchor=north},legend columns=3]
\addplot[omThm,thick,mark=*,mark size=1.3pt] table[x=k,y=mean20]{figdata/hft/13-index-and-basket-arbitrage/edge.csv};
\addplot[omDef,thick,dashed,mark=square*,mark size=1.3pt] table[x=k,y=mean2]{figdata/hft/13-index-and-basket-arbitrage/edge.csv};
\addplot[gray,thin,mark=triangle*,mark size=1.3pt] table[x=k,y=costs]{figdata/hft/13-index-and-basket-arbitrage/edge.csv};
\legend{{bord net, une jambe toutes les 20 $\mu$s},{bord net, une jambe toutes les 2 $\mu$s},{coûts}}
\end{axis}
\end{tikzpicture}

\omcaption{Edge net par transaction (capture moins coûts, en points de base du notionnel) en fonction de la taille du \omterm{def:hf:index-and-basket-arbitrage:partial}{panier partiel}, lorsque le contrat à terme bondit de 3
points de base et que la cotation de chaque action suit après un délai exponentiel de moyenne 400 microsecondes~; jambes envoyées toutes les 20 ou toutes les 2 microsecondes~; coûts
sont les demi-fourchettes des jambes et du contrat à terme à l'entrée et à la sortie. 20\,000 transactions par point. Données~: \texttt{hf\_basket.sweep}.}\label{fig:hf:index-and-basket-arbitrage:edge}
\end{omfigure}

\section{Financement et risque de dividende}

Détenue au-delà de la journée, la position porte les termes du \omterm{def:hf:fair-value:fair}{prix juste}~: le taux de financement (le taux de financement implicite du livre~1, chapitre~21, comparé à
celui de l'arbitragiste) et les dividendes. Un rendement de dividende prévu avec une erreur de 0.2~\% par an déplace le \omterm{def:hf:fair-value:fair}{prix juste} d'un contrat à terme à trois mois de 5 points de base
soit davantage que la totalité de l'edge de haute fréquence~; le risque de dividende (livre~5, chapitre~5) explique pourquoi la version haute fréquence du trade est clôturée
en quelques secondes, et la version lente constitue un ouvrage à part, avec ses propres prévisions de dividendes.

\section{Le même indice sur plusieurs contrats à terme}

\begin{definition}[Contrats à terme multicotés]\label{def:hf:index-and-basket-arbitrage:cross}
Les \emph{contrats à terme multicotés}\index{contrats à terme multicotés} sont des contrats à terme sur le même indice, cotés sur des bourses différentes ou dans des tailles ou devises différentes~;
leurs prix sont liés par l'indice, et leurs écarts reflètent les différences de devise, de conditions des contrats, d'heures de négociation et de liquidité.
\end{definition}

Le même indice se négocie sous forme de contrat de pleine taille, de contrat mini, de contrat micro, et sur des bourses situées dans plusieurs fuseaux horaires. Une position sur l'un se couvre sur l'autre dans le rapport de leurs multiplicateurs et de leurs devises~: un contrat d'Osaka de
\textyen1\,000 fois le Nikkei est couvert par $1\,000/(5\times150)=1.33$ contrats d'une cotation hypothétique libellée en dollars
rapportant 5~\$ par point, à 150 yens pour un dollar. L'arbitrage entre les deux est un trade lead--lag (chapitre~8) entre deux carnets d'ordres portant sur
un même indice, avec la devise comme troisième jambe lorsque les contrats sont libellés différemment, et un ajustement quanto dans l'écart juste.

\begin{dated}{2026-09}{Contrats micro et un même indice sur trois bourses}\label{dat:hf:index-and-basket-arbitrage:listings}
CME Group a lancé des contrats à terme Micro E-mini sur les indices S\&P 500, Nasdaq-100, Russell 2000 et Dow Jones Industrial Average le 6 mai 2019, d'une taille égale au dixième
de celle de ses contrats E-mini. Les contrats à terme sur le Nikkei 225 se négocient au Singapore Exchange et chez CME Group (\textyen500 fois l'indice, en yens ou en dollars,
tick de 5 points) et à l'Osaka Exchange (\textyen1\,000 fois l'indice, en yens, tick de 10 points), comme le décrit une comparaison de courtier de juin 2024~;
dans leur système de compensation mutuelle, une position ouverte chez CME ou SGX peut être clôturée sur l'autre le même jour.
\end{dated}

\section{Fiches de stratégie}

\begin{strategyfile}{Contrat à terme contre panier partiel optimisé}\label{strat:hf:index-and-basket-arbitrage:partial}
\sfield{Qui paie, et pourquoi} Les teneurs de marché des composantes dont les cotations accusent un retard sur un mouvement de l'ensemble du marché apparu d'abord sur le contrat à terme.
\sfield{Instruments et plateformes} Un contrat à terme sur indice et un \omterm{def:hf:index-and-basket-arbitrage:partial}{panier partiel} de ses composantes, négociés sur les plateformes transparentes.
\sfield{Signal} Le mouvement du contrat à terme au-delà du \omterm{def:hf:fair-value:fair}{prix juste}, converti en points d'indice, comparé à la valeur cotée du panier.
\sfield{Dimensionnement et exécution} Acheter (vendre) le \omterm{def:hf:index-and-basket-arbitrage:partial}{panier partiel} jambe par jambe, les plus fortes pondérations en premier, et vendre (acheter) le contrat à terme~; choisir la taille du panier
d'après la vitesse de la passerelle et les retards de cotation des composantes~; débouclage en quelques secondes, à mesure que les actions rattrapent leur retard.
\sfield{Coûts} Demi-fourchettes sur chaque jambe à l'entrée et à la sortie (plus larges pour les petites valeurs), celles du contrat à terme, les frais~; tracking error pendant la détention.
\sfield{Comment elle s'éteint} Vitesse~: avec des jambes toutes les 2 microsecondes au lieu de 20, le meilleur panier passe de 5 noms à 10 et l'edge est divisé
par deux~; l'arbitragiste doté de la passerelle la plus rapide fixe le prix de chacun des autres, plus lents.
\sfield{Horizon, capacité, infrastructure} De la microseconde à la seconde~; capacité tirée des tailles au meilleur prix des composantes~; une passerelle qui envoie de nombreux
ordres à la fois (livre~13).
\sfield{Backtest rigoureux} Instants de mise à jour propres à la cotation de chaque action après le mouvement du contrat à terme~; instants d'envoi réels des jambes~; exécutions seulement aux cotations
qui existaient à l'arrivée de l'ordre.
\sfield{Sources} Roll, Schwartz et Subrahmanyam (2007)~; ce chapitre.
\end{strategyfile}

\begin{strategyfile}{Arbitrage de contrats à terme sur indice multicotés}\label{strat:hf:index-and-basket-arbitrage:cross}
\sfield{Qui paie, et pourquoi} Les acteurs de la cotation la moins liquide, dont les prix suivent tard ceux de la cotation qui mène.
\sfield{Instruments et plateformes} Les contrats à terme sur un même indice cotés sur deux bourses (le Nikkei à Osaka, Singapour et Chicago), et la devise.
\sfield{Signal} L'écart entre les cotations, au-delà de son \omterm{def:hf:fair-value:fair}{prix juste} (devise, ajustement quanto, financement, différences d'échéance et de règlement-livraison).
\sfield{Dimensionnement et exécution} Prendre la cotation en retard, couvrir sur celle qui mène dans le rapport des multiplicateurs et des devises~; conserver jusqu'à la convergence ou compenser
entre bourses lorsqu'un système de compensation mutuelle le permet.
\sfield{Coûts} Les fourchettes et frais des deux contrats à terme, la couverture de devise, les appels de marge auprès de deux chambres de compensation.
\sfield{Comment elle s'éteint} Le leader change selon l'heure de la journée~: chaque cotation mène pendant les heures de son propre marché.
\sfield{Horizon, capacité, infrastructure} Des millisecondes (en colocation) aux minutes~; liaisons à longue distance entre les centres de données des bourses.
\sfield{Backtest rigoureux} Horloges synchronisées entre continents~; cours de la devise à chaque transaction~; données du calcul de l'écart juste telles qu'elles étaient connues à ce moment.
\sfield{Sources} Conditions des contrats (encadré daté)~; chapitre~8.
\end{strategyfile}

\begin{strategyfile}{Arbitrage entre contrat normal et contrat mini}\label{strat:hf:index-and-basket-arbitrage:mini}
\sfield{Qui paie, et pourquoi} Les traders du contrat le plus petit, dont le carnet d'ordres se met à jour après celui du plus grand.
\sfield{Instruments et plateformes} Un contrat normal ou E-mini et son contrat mini ou micro sur la même bourse.
\sfield{Signal} L'écart de cours entre les deux, en points d'indice (leur écart juste est nul pour la même échéance et le même règlement-livraison).
\sfield{Dimensionnement et exécution} Dix micros contre un E-mini~; coter le micro autour du cours du E-mini, ou prendre la transaction lorsque l'écart dépasse les coûts.
\sfield{Coûts} Les frais par contrat, qui pèsent dix fois plus sur le notionnel du micro~; le tick est le même en points d'indice.
\sfield{Comment elle s'éteint} Beaucoup d'intervenants cotent les deux carnets à partir des mêmes données~; l'écart dépasse rarement un tick.
\sfield{Horizon, capacité, infrastructure} Des microsecondes~; capacité déterminée par la profondeur du petit contrat.
\sfield{Backtest rigoureux} L'appariement des ordres par la bourse entre les deux carnets~; frais par contrat.
\sfield{Sources} Lancement du Micro E-mini par le CME Group (encadré daté).
\end{strategyfile}

\section{Tutoriel~: trente noms parmi cinq cents}\label{tut:hf:index-and-basket-arbitrage:tut}

\begin{tutorial}
\textbf{Objectif.} Trouver la taille de panier qui maximise l'edge net d'un arbitrage d'indice à haute fréquence, et sa dépendance à la vitesse des
jambes. \textbf{État final~:} la \cref{fig:hf:index-and-basket-arbitrage:edge} et le tableau ci-dessous.
\begin{tutsteps}
\item \textbf{L'indice.} \texttt{firm.basketarb.Index}~: cinquante noms avec des poids de capitalisation lognormaux, des bêtas compris entre 0.6 et 1.4 mis à l'échelle pour que
celui de l'indice vaille un, une volatilité idiosyncratique de 1 à 3~\% par jour, des demi-fourchettes allant de 0.5 point de base pour le plus grand nom à 1.7 pour le plus petit.
\item \textbf{Le \omterm{def:hf:index-and-basket-arbitrage:partial}{panier partiel}} (\cref{def:hf:index-and-basket-arbitrage:partial}), choisi de façon gloutonne.
\omcode{code/firm/basketarb/firm_basketarb.py}{39}{61}{Choix glouton des noms~; poids des moindres carrés sur l'ensemble retenu.}\label{lst:hf:index-and-basket-arbitrage:partial}
\item \textbf{La transaction.} Le future saute de 3 points de base~; le cours de chaque action suit après un délai exponentiel de moyenne 400 microsecondes~; les
jambes partent toutes les 20 (ou 2) microsecondes, la plus grande en premier~; on conserve la position dix secondes.
\omcode{code/firm/basketarb/firm_basketarb.py}{97}{110}{Les jambes dans le temps captent le mouvement de leur action~; coûts à l'entrée et à la sortie~; bruit de suivi pendant la détention.}\label{lst:hf:index-and-basket-arbitrage:trade}
\item \textbf{Balayage} de la taille du panier, d'un nom à cinquante (\texttt{hf\_basket.sweep}).
\end{tutsteps}
\textbf{Pour aller plus loin.} Envoyer les jambes en parallèle sur plusieurs passerelles~; faire dépendre les délais de cotation de la liquidité de chaque action~; ajouter
l'avance propre du future sur un second future, plus lent.
\end{tutorial}

\begin{center}
\small
\begin{tabular}{@{}lrrrrrr@{}}
\toprule
noms & 1 & 5 & 10 & 15 & 30 & 50\\
\midrule
capture (pb) & 1.44 & 2.13 & 2.17 & 2.10 & 1.84 & 1.72\\
coûts (pb) & 0.93 & 1.21 & 1.29 & 1.39 & 1.52 & 1.65\\
écart-type du suivi (pb) & 1.74 & 0.92 & 0.62 & 0.44 & 0.16 & 0\\
edge net (pb) & 0.51 & 0.92 & 0.88 & 0.71 & 0.32 & 0.07\\
edge / écart-type & 0.29 & 0.91 & 1.24 & 1.33 & 1.06 & 0.32\\
edge net, jambes toutes les 2 $\mu$s (pb) & 0.57 & 1.23 & 1.36 & 1.34 & 1.28 & 1.15\\
\bottomrule
\end{tabular}
\end{center}

Avec une jambe toutes les vingt microsecondes, l'edge net culmine à cinq noms (0.92 point de base par transaction) et l'edge par unité de risque à quinze. Acheter
les cinquante noms capte moins (les dernières jambes arrivent une milliseconde après le mouvement du future, quand la plupart des cotations ont déjà été mises à jour) et coûte davantage (les fourchettes des petits
noms) pour une transaction presque sans risque et presque sans valeur~: 0.07 point de base. Des jambes dix fois plus rapides déplacent l'optimum de l'edge vers dix noms (1.36
points de base) et font du panier complet le meilleur rapport par unité de risque. Lorsque les cotations des actions suivent plus vite (une moyenne de 150 microsecondes), le meilleur panier
compte deux noms~; plus lentement (une milliseconde), dix.

\section{Construction~: le module d'arbitrage de panier}\label{bld:hf:index-and-basket-arbitrage:basketarb}

\begin{build}
\bfield{Rôle} Choisir des \omterm{def:hf:index-and-basket-arbitrage:partial}{paniers partiels}, fixer des bandes de base, simuler le legging et dimensionner les couvertures sur \omterm{def:hf:index-and-basket-arbitrage:cross}{contrats à terme multicotés}.
\bfield{Interface} \texttt{covariance}, \texttt{tracking\_var}, \texttt{partial\_basket(w, cov, k)}, \texttt{fair\_future}, \texttt{basis\_band},
\texttt{Index(n, seed)}, \texttt{trade(index, k, jump\_bp, leg\_us, lag\_us, hold\_s)}, \texttt{cross\_listed}, \texttt{hedge\_ratio}.
\bfield{Règles} Les poids sont des moindres carrés sur l'ensemble retenu~; les jambes partent de la plus grande à la plus petite~; tous les P\&L sont exprimés en points de base du notionnel.
\bfield{Tests d'acceptation} \texttt{code/firm/basketarb/tests/}~: covariance et variance de suivi calculées à la main~; la variance de suivi diminue à chaque nom et
vaut zéro pour l'indice complet~; le couple glouton n'est pas meilleur que le meilleur couple obtenu par force brute~; future juste et bande calculés à la main~; des jambes plus rapides captent davantage à
coût égal~; avec toutes les jambes dans le temps, la capture est le saut.
\bfield{Pour aller plus loin} Passerelles en parallèle~; pas de cotation et lots minimaux~; ordonnancement optimal des jambes selon la vitesse de mise à jour des cotations plutôt que selon le poids.
\end{build}

\begin{omsources}
\item R.~Roll, E.~Schwartz, A.~Subrahmanyam, Liquidity and the law of one price: the case of the futures-cash basis, \emph{Journal of Finance}
62(5), 2007, 2201--2234.
\item CME Group, annonce du lancement des contrats à terme micro sur indice actions E-mini, 6 mai 2019.
\item Phillip Nova, Differences between the Nikkei 225 futures contracts on SGX, CME Group and OSE, 5 juin 2024.
\end{omsources}

\section{Exercices}

\begin{exercise}[$\star$]\label{exo:hf:index-and-basket-arbitrage:1}
Un indice vaut 5\,000, le taux est de 5~\%, le rendement du dividende de 1.5~\%, et le future expire dans trois mois. Quelle est sa juste valeur, et la bande à 1
point de base de coûts~?
\end{exercise}

\begin{exercise}[$\star$]\label{exo:hf:index-and-basket-arbitrage:2}
Combien de contrats Micro E-mini couvrent un E-mini~? Combien de contrats d'une hypothétique cotation du Nikkei libellée en dollars, rapportant 5~\$ par point, couvrent un contrat d'Osaka de
\textyen1\,000 par point, à 150 yens pour un dollar~?
\end{exercise}

\begin{exercise}[$\star$]\label{exo:hf:index-and-basket-arbitrage:3}
De combien de points de base une prévision de dividende fausse de 0.2~\% par an modifie-t-elle la juste valeur d'un future à trois mois~?
\end{exercise}

\begin{exercise}[$\star\star$]\label{exo:hf:index-and-basket-arbitrage:4}
Avec des jambes toutes les 20 microsecondes et des délais de cotation de moyenne 400, quelle est la probabilité que la dixième jambe arrive à temps~? La trentième~?
\end{exercise}

\begin{exercise}[$\star\star$]\label{exo:hf:index-and-basket-arbitrage:5}
Pourquoi l'edge par unité de risque culmine-t-il à un panier plus grand que l'edge lui-même~?
\end{exercise}

\begin{exercise}[$\star\star$]\label{exo:hf:index-and-basket-arbitrage:6}
Pourquoi envoie-t-on d'abord les plus grands noms~?
\end{exercise}

\begin{exercise}[$\star\star\star$]\label{exo:hf:index-and-basket-arbitrage:7}
\emph{Programmation.} Relancer le balayage avec des délais de cotation de moyenne 150 et 1\,000 microsecondes. Quelles tailles de panier maximisent l'edge et l'edge par unité de
risque~?
\end{exercise}

\begin{exercise}[$\star\star\star$]\label{exo:hf:index-and-basket-arbitrage:8}
\emph{Trouver la faille.} «~Notre backtest achète les 500 noms aux cours affichés au moment où le future a bougé~; le panier complet n'a aucune erreur de suivi, c'est donc
la meilleure transaction.~»
\end{exercise}

\section{Problème~: trente noms sur cinq cents}

\begin{problem}[{Problème du week-end --- trente noms sur cinq cents}]\label{pb:hf:index-and-basket-arbitrage:1}
Un arbitragiste négocie un future sur indice contre ses composantes lorsque le future bouge le premier.

\medskip\noindent\textbf{Partie I --- La base.}
\begin{enumerate}
\item Écrire le fair value et la bande d'absence d'arbitrage.
\item Pourquoi la base sort-elle de sa bande à haute fréquence, et pourquoi y revient-elle~?
\item Qu'ont établi Roll, Schwartz et Subrahmanyam sur la liquidité et la base~?
\item Pourquoi les prévisions de dividendes importent-elles pour la transaction lente et non pour la rapide~?
\end{enumerate}

\medskip\noindent\textbf{Partie II --- Paniers et jambes.}
\begin{enumerate}[resume]
\item Définir un \omterm{def:hf:index-and-basket-arbitrage:partial}{panier partiel} et sa variance de suivi.
\item Comment sont choisis les noms et les pondérations~?
\item Définir le \omterm{def:hf:index-and-basket-arbitrage:legging}{risque d'exécution par jambes} et écrire la probabilité que la jambe $j$ soit exécutée à temps.
\item Pourquoi chaque nom ajouté coûte-t-il plus que le précédent~?
\end{enumerate}

\medskip\noindent\textbf{Partie III --- Mesures.}
\begin{enumerate}[resume]
\item Donner la capture, les coûts et l'edge net pour 1, 5, 15 et 50 noms avec 20 microsecondes par jambe.
\item Où l'edge et l'edge par unité de risque atteignent-ils leur maximum~?
\item Qu'est-ce qui change avec des jambes toutes les 2 microsecondes~?
\item Qu'est-ce qui change lorsque les cotations suivent en 150 microsecondes ou en une milliseconde~?
\end{enumerate}

\medskip\noindent\textbf{Partie IV --- Le verdict.}
\begin{enumerate}[resume]
\item Énoncer le \emph{résultat nommé}~: l'edge net par transaction en fonction de la taille du panier, et la taille qui le maximise.
\item Définir les \omterm{def:hf:index-and-basket-arbitrage:cross}{contrats à terme multicotés} et le ratio de couverture entre deux d'entre eux.
\item Que dit l'encadré daté sur les contrats micro et les cotations du Nikkei~?
\item Laquelle des trois stratégies dépend le plus de la vitesse~?
\item Pourquoi l'arbitrage améliore-t-il le marché pour tous les autres~?
\item Que changerait un indice à 500 noms~?
\item Comment estimer les délais de cotation des actions à partir des données~?
\item En une phrase~: qu'est-ce qui détermine la taille du panier~?
\end{enumerate}
\end{problem}

\section{Questions d'entretien}

\begin{interviewq}[$\star$ \iqroles{trader}]\label{iq:hf:index-and-basket-arbitrage:1}
Le E-mini se situe 2 points au-dessus de son fair value. Décrivez ce que vous vérifiez avant de trader.
\end{interviewq}

\begin{interviewq}[$\star\star$ \iqroles{chercheur}]\label{iq:hf:index-and-basket-arbitrage:2}
Dérivez les pondérations des moindres carrés d'un \omterm{def:hf:index-and-basket-arbitrage:partial}{panier partiel} et montrez qu'elles minimisent la variance de suivi pour les noms retenus.
\end{interviewq}

\begin{interviewq}[$\star\star$ \iqroles{développeur}]\label{iq:hf:index-and-basket-arbitrage:3}
Vous devez envoyer 50 ordres aussi vite que possible après un paquet de données de marché. Qu'est-ce qui vous limite, et comment les faites-vous sortir en parallèle~?
\end{interviewq}

\begin{interviewq}[$\star\star$ \iqroles{risque}]\label{iq:hf:index-and-basket-arbitrage:4}
La moitié des jambes d'une transaction de panier ont été exécutées et les autres non. Quelle est votre exposition, et que doit faire le système~?
\end{interviewq}

\begin{interviewq}[$\star\star$ \iqroles{trader}]\label{iq:hf:index-and-basket-arbitrage:5}
Le Nikkei se négocie à Osaka, à Singapour et à Chicago. Lequel mène, et quand~?
\end{interviewq}

\begin{interviewq}[$\star\star\star$ \iqroles{chercheur}]\label{iq:hf:index-and-basket-arbitrage:6}
Avec un espacement des jambes $\delta$, des délais de cotation exponentiels de moyenne $\tau$ et des pondérations égales, déterminer le nombre de jambes
qui maximise la capture espérée diminuée d'un coût par jambe.
\end{interviewq}
